\documentclass[11pt]{article}
\usepackage[margin=1in]{geometry}
\usepackage{amsmath,amssymb,amsfonts,amsthm,mathtools,bm}
\usepackage[T1]{fontenc}
\usepackage{microtype}
\newtheorem{theorem}{Theorem}
\newtheorem{lemma}{Lemma}
\newtheorem{cor}{Corollary}
\newtheorem{defn}{Definition}
\newtheorem{principle}{Principle}
\newtheorem{statement}{Statement}
\theoremstyle{remark}
\newtheorem{statementr}{Statement}

\title{Two Algorithms for Exact Quantum Dynamics from Trajectories}
\author{James Taylor}
\date{Working draft derived from the 2003 thesis}
\begin{document}
\maketitle

\section{Two 
algorithms for exact quantum dynamics from trajectories} \label{c:algo}


In this chapter, we shall discuss a new idea for 
solving Schr\"odinger's equation. The results are few, but the potential for more is
extremely high. 
We start by rewriting Schr\"odinger's equation into the usual pair of
the modified 
Hamilton-Jacobi equation and the continuity equation. Then a
natural iterative process suggests itself. Namely, ignore the quantum
potential, solve the H-J equation, use the continuity equation to get a new
density, use that density to obtain an approximate quantum potential, and
repeat. The fixed point of this scheme is the solution. We are
particularly interested in an equivalent 
formulation involving trajectories.
 An important fact is that if we had the Bohmian\
trajectories and the initial wave function, then we could produce the solution
to Schr\"odinger's equation. But to find the Bohmian\ trajectories,  we need to already have
the solution to Schr\"odinger's equation. The two
algorithms we will be discussing provide a separation. We start with
an initial set of trajectories, using them to construct an approximate
wave function. We then use that new wave function\ to find new trajectories. We
continue. Pictorially, we have 
\[
Q_0 \to \ensuremath{\psi}_1 \to Q_1 \to \cdots \to Q_n \to \ensuremath{\psi}_{n+1} \to Q_{n+1}
\to \cdots 
\]
Convergence is a central question of any such proposal. Unfortunately,
we have no general answer to that question. We do, however, have some examples
of beautiful convergence. Unfortunately, these are restricted to only
quadratic
potentials in 1 dimension. Additionally, we have used the process on
the fixed point to derive a PDE that the family of Bohmian
trajectories satisfies. That is to say, we can actually find the
Bohmian trajectories before solving Schr\"odinger's equation. We explicitly do this in the
special case of when the Bohmian trajectories are linear in the
initial configuration; we
do this only in 1 dimension. 

\section{Background material}

In chapter 3, we derived the Hamilton-Jacobi equation and the
continuity equation from Schr\"odinger's equation. In this section, we write down what the
solutions are to these equations in terms of the characteristics. 
It is convenient to start with some
definitions and notations. Appendix \ref{a:app7} contains the proofs
of these assertions as well as a discussion of the derivatives of the
determinant of the differential of a map. 

We start with some definitions. 
By \emph{family of trajectories}, we shall mean
essentially a flow although we do not make assumptions about
invertibility. If $Q$ is a family of trajectories, then we shall write
$Q(x_0, t)$  to denote the configuration at time $t$ of a system whose
initial configuration was $x_0$. We also define the action
functional. Given a family of trajectories $Q$, a potential $U$, and 
an initial function $T$, we define 


\[ 
\ensuremath{\mathcal{S}}[Q, U, T] (x_0, t) \ensuremath{:=}   T(x_0) + \int_0^t [
\frac{\dot{Q}^2}{2}(x_0, s) -
U(Q(x_0,s), s)] ds.
\]
Our last definition is for the amplitude functional. Given the family
$Q$ and an initial function $P$, we define 
\[
\ensuremath{\mathcal{R}}[Q, P](x_0, t) = P(x_0) |DQ(x_0,t) |^{-\frac{1}{2}}. 
\]
As we shall discuss, these are essentially the solutions to the
Hamilton-Jacobi and continuity equations, respectively. 

In this chapter, all wave functions will be complex-valued wave function
s over a given manifold \ensuremath{\mathcal{Q}}; the examples will all be for the
simple manifold 
$\ensuremath{\mathcal{Q}} = \ensuremath{\mathbb{R}^{}}$. 
We generally demand these functions to be smooth elements of
$L^2(\ensuremath{\mathcal{Q}})$. We start 
by writing the wave function, \ensuremath{\psi}, in polar form;   
$
\ensuremath{\psi} = Re^{\frac{iS}{\ensuremath{\hbar}}}, 
$
where the amplitude $R \geq 0$, and the phase, $S$, is real. The
quantum density, \ensuremath{\rho}, is $\ensuremath{|\ensuremath{\psi}|^2}  = R^2.$ 
We first think about the setup for $N$ particles moving in
$\ensuremath{\mathbb{R}^{k}}$; we shall use the usual setup for this rather than the one
we discussed in Chapter \ref{c:manpart}. 
Solving Schr\"odinger's equation\ is
equivalent to solving the coupled equations  


\begin{align*}
\ensuremath{\frac{\partial {S}}{\partial {t}}} &= -\sum_{k=1}^N \frac{1}{2m_k}(\nabla_k S \cdot \nabla_k
S) - {V} + \sum_{k=1}^N \frac{\ensuremath{\hbar}^2}{2m_k} \frac{\ensuremath{\Delta}_k R}{R} 
\\
\ensuremath{\frac{\partial {\ensuremath{\rho}}}{\partial {t}}} &= - \sum_{k=1}^N  \nabla_k \cdot (\ensuremath{\rho} \frac{1}{m_k}\nabla_k S)
.\\
\end{align*}

For a general manifold, one immediately has the same thing using the
Riemannian metric. One can also formulate $N$ particles moving in
\ensuremath{\mathcal{Q}}. The masses can be absorbed into the metric as described in
Chapter \ref{c:geom}. Thus, we can drop the masses and the
summations. We could also write the above equations explicitly using
the metric: 
\begin{align*}
\ensuremath{\frac{\partial {S}}{\partial {t}}} &= -\frac{1}{2}G(DS, DS)  - {V} + \frac{\ensuremath{\hbar}^2}{2}
\frac{G\circ D^2 R}{R} = - \frac{1}{2} \nabla S \cdot \nabla S - V + 
\tag{H-J}  \frac{\ensuremath{\hbar}^2}{2}
\frac{\ensuremath{\Delta} R}{R}
\\
\ensuremath{\frac{\partial {\ensuremath{\rho}}}{\partial {t}}} &= - G\circ(D(\ensuremath{\rho} DS)) = -\nabla \cdot (\ensuremath{\rho} \nabla S).
 \tag{C}
\\
\end{align*}
We shall absorb the mass into the metric except possibly when we do
the examples. 


Equation (H-J) is almost the Hamilton-Jacobi equation for classical motion
under the potential $V$; it differs by the  
 quantum potential,  
 $ -\frac{\ensuremath{\hbar}^2}{2}\frac{\ensuremath{\Delta} R}{R} $.
Classically, the velocity field is $\nabla S$, which is just the
Bohmian\
velocity field. Equation  
(C) is the continuity equation; this describes how \ensuremath{\rho}\
evolves under the flow corresponding to configurations
moving with the Bohmian\ velocities.  The trajectory of the system is a single
integral curve of the Bohmian\ 
velocity field.  

Let $Q(x_0,t)$ be the family of
trajectories solving $\dot{Q} (x_0,t)  = \nabla S (Q(x_0,t))$ 
with inverse $Q^{-1}$ and  $\ensuremath{\lvert {DQ^{-1}} \rvert} \neq 0$.  Let $\ensuremath{\rho}_0$ and
$S_0$ be the amplitude and phase of the initial wave function. Then in the
appendix, we show that $(\ensuremath{\mathcal{R}}[Q, \ensuremath{\rho}_0] (Q^{-1})(x,t)))^2$ is the solution
to (C) while $\ensuremath{\mathcal{S}}[Q, V -\frac{\ensuremath{\hbar}^2}{2}\frac{\ensuremath{\Delta} R}{R}, S_0]
(Q^{-1}(x,t))$ is the solution to (H-J). 


\section{The algorithms}


We start by describing the algorithms in terms of the approximating
PDEs.
We
shall begin with algorithm A1. At the beginning
of the \ensuremath{\text{n}^{\text{th}}}\ step, we have 
$\ensuremath{\rho}_{n-1} \ensuremath{:=} R_{n-1}^2$ and $S_{n-1}$. We demand that $S_n$ satisfies
\[ \ensuremath{\frac{\partial {S_n}}{\partial {t}}} =  -\nabla S_n \cdot \nabla S_{n-1} +
\frac{1}{2} \nabla S_{n-1} \cdot \nabla S_{n-1} - V  +
\frac{\ensuremath{\hbar}^2 \ensuremath{\Delta} R_{n-1}}{2 R_{n-1}} \tag{H-J:A1, $n$}
\]
and that $\ensuremath{\rho}_{n}$ satisfies 
\[
\ensuremath{\frac{\partial {\ensuremath{\rho}_n}}{\partial {t}}} = - \nabla \cdot (\ensuremath{\rho}_n \nabla S_{n}) \tag{C:A1, $n$}.
\]
For algorithm A2, we have a similar setup, but we demand that $S_n$
satisfies 
\[
 \ensuremath{\frac{\partial {S_n}}{\partial {t}}} = -\frac{1}{2} \nabla S_{n} \cdot \nabla S_{n} - V  +
\frac{\ensuremath{\hbar}^2 \ensuremath{\Delta} R_{n-1}}{2 R_{n-1}} \tag{H-J:A2, $n$}
\]
and that $\ensuremath{\rho}_{n}$ satisfies 
\[
\ensuremath{\frac{\partial {\ensuremath{\rho}_n}}{\partial {t}}} = - \nabla \cdot (\ensuremath{\rho}_n \nabla S_{n}) \tag{C:A2, $n$}.
\]
We denote by $U_{n-1}$  the potential $V  - \frac{\ensuremath{\hbar}^2}{2}\frac{\ensuremath{\Delta}
R_{n-1}}{ R_{n-1}}$. 

There are several comments that may be made. First, the fixed point of
both algorithms provides the solution of Schr\"odinger's equation. Second, both equations in
A1 are linear, understood as solved in the appropriate order, 
 whereas  (H-J:A2) is quadratic in the spatial
derivative. Third, in the limit 
$\ensuremath{\hbar} \to 0$, A2 becomes classical mechanics whereas A1 does
not. Indeed, if $\ensuremath{\hbar} =0$, A2 does not iterate and the only step becomes a
quasiclassical approximation. But for A1, this is not the case. In
fact, if $\ensuremath{\hbar} =0$, then A1 still iterates. It becomes an algorithm for
computing the solution to the classical H-J equation. 


On the \ensuremath{\text{n}^{\text{th}}}\ step of either algorithm, we define $\dot{Q}_n \ensuremath{:=}
\nabla S_n$. With this definition, $Q_n$  corresponds to the characteristics
of (H-J:A1, 
$n+1$), (C:A1, $n$), (H-J:A2,
$n$) and (C:A2, $n$); \emph{characteristics} are trajectories that
convert a PDE to an ODE. Indeed, in A1, it can be shown that 
\[ S_{n} (x,t) = \ensuremath{\mathcal{S}}[Q_{n-1}, U_{n-1}, S(\cdot, 0)] (Q_{n-1}^{-1} (x, t), t) 
\]
 and 
\[R_n = \ensuremath{\mathcal{R}}[Q_n, R(\cdot, 0)] (Q_n^{-1} (x,t), t).
\]
In A2, we have 
\[ S_{n} (x,t) = \ensuremath{\mathcal{S}}[Q_{n}, U_{n-1}, S(\cdot, 0)] (Q_n^{-1} (x, t), t) 
\]and \[
R_n = \ensuremath{\mathcal{R}}[Q_n, R(\cdot, 0)] (Q_n^{-1} (x,t), t).
\]
From this perspective, A1 is close to Bohmian mechanics\ in spirit. At every stage,
we start with $S_n$ and compute the integral curves of its gradient
thereby forming the family $Q_n(x_0,t)$. 
Then  we compute the Jacobian of $Q_n$ to obtain $R_{n+1}$ and
the action integral using the new potential $U_{n+1}$to obtain $S_{n+1}$. In
contrast, A2 is 
closer to the spirit of classical mechanics. To find $Q_n$,
we actually solve the classical equations with the force being
$-\nabla U_{n-1}$. We then use the Jacobian of $Q_n$ to define $R_n$ and
$U_n$. This sets us up to calculate the next set of trajectories. In A2, we
do not need to compute $S_n$ in order to complete the step. 


In algorithm A2, we encounter the issue that $S_n$ may not be
well-defined for all time; this generally corresponds to $Q_n$ not being
invertible. Following \cite[pages~438--445]{Arn89},  we define
 an approximate wave function\ by 
\[
{\ensuremath{\psi}}_n (x, t)  \ensuremath{:=} \sum_{ \{y_j | Q_n(y_j, t) = x \} } \ensuremath{\mathcal{R}} [Q_n,
 R(\cdot,0)] (y_j,  t)  
e^{i(\ensuremath{\mathcal{S}}[Q_n, U_{n-1}, S(\cdot, 0)] (y_j, t) -
 \frac{\pi}{2} \mu_j)}.
\]
We then obtain $R_n$ and $S_n$ from the polar form of that wave function. 
 The Maslov
index $\mu_j$ gives us an appropriate phase for the summation. 
 Appendix  \ref{a:app7} gives a slighly more informative discussion of this. 

\subsection{The algorithms in short} \label{as:short}

We now describe the algorithms in a more mechanical fashion. Given $T(x)$,
$P(x)$, $V(x,t)$, and $Q(x,t)$, the two algorithms may be viewed as generating
different ODEs to solve involving the function $q(x,t)$ once the
initial $x$ is fixed. 

We shall denote the following procedure as $A1(T, P, V, Q)$; it
produces a family of trajectories $q(x, t)$.
\begin{enumerate}
\item Compute $H \ensuremath{:=} Q^{-1}$. 
\item Define $R \ensuremath{:=} P(H) \ensuremath{\lvert {DQ^{-1}} \rvert}^{\frac{1}{2}}$.
\item Compute $\tilde{S} \ensuremath{:=} T + \int_0^t [\frac{\dot{Q}^2}{2} -
V(Q,s) -  \frac{\ensuremath{\hbar}^2}{2}\frac{\ensuremath{\Delta} R}{R} (Q,s)] ds$.
\item Define $S \ensuremath{:=} \tilde{S} (H)$.
\item Define $q(x,t) \ensuremath{:=} \gamma_x(t)$ where $\gamma_x$ the solution of
$\dot{\gamma}_x (s) = (\nabla S)(\gamma_x(s), s)$ with initial condition
$\gamma_x(0) = x$ .  
\end{enumerate}
For an algorithm, fix $T$, $P$, $V$ and define $B(Q) \ensuremath{:=}
A1(T,P,V,Q)$. 
The \ensuremath{{n}^{th}}\ step of A1 is then the same as $B^n(Q_0)$. To start, one
can choose a best guess, $Q_0(x,t) = x$, or the solution to 
$\dot{Q}_0(x,t) =
\nabla S_0(Q)$.

We do a similar iteration for A2. We first describe the version when
$Q$ is invertible. 
We shall denote the following procedure as $A2(T, P, V, Q)$; it
produces a family of trajectories $q(x, t)$.
\begin{enumerate}
\item Compute $H \ensuremath{:=} Q^{-1}$. 
\item Define $R \ensuremath{:=} P(H) \ensuremath{\lvert {DQ^{-1}} \rvert}^{\frac{1}{2}}$.


\item Define $q(x,t) \ensuremath{:=} \gamma_x(t)$ where $\gamma_x$ the solution of
$\ddot{\gamma}_x (s) = -(\nabla (V  -  \frac{\ensuremath{\hbar}^2}{2m}\frac{\ensuremath{\Delta}
R}{R}) )(\gamma_x(s), s)$ with initial condition 
$\gamma_x(0) = x$, $\dot{\gamma}_x(0) = \nabla T (x)$.  
\end{enumerate}
For an algorithm, fix $T$, $P$, $V$ and define $C(Q) \ensuremath{:=}
A2(T,P,V,Q)$. 
The \ensuremath{{n}^{th}}\ step of A1 is then the same as $C^n(Q_0)$. To start, one
can choose a best guess, $Q_0(x,t) = x$, the solution to $\dot{Q}_0(x,t)=
\nabla S_0(Q)$, or the classical trajectories using the initial $S$ to
for the initial conditions. 


Finally, we have the case when the trajectories are not invertible, but are
generated in a classical way. In particular, we shall let $U$ be the
potential which generated the trajectories. In the above algorithm,
it was not necessary to have this information.
As the above is mostly a special case of what is below, 
we shall denote the following procedure  as $A2(T, P, V, Q,U)$; it
produces a family of trajectories $q(x, t)$ and the potential $u$ that
generated those trajectories.
\begin{enumerate}
\item Define the set of inverse points $H_{y,t} \ensuremath{:=} Q^{-1} (y,t)$. 
\item Define $\mu(x,t)$ to be the Maslov index of the curve $Q (x,t)$
under the flow of $Q$. In particular, it is the number of signed zeroes that
the trajectory has crossed of the function $\ensuremath{\lvert {DQ} \rvert} (x, t)$.
\item Compute $S (x,t) \ensuremath{:=} T(x) + \int_0^t [\frac{\dot{Q_x}^2}{2}(s) 
- U(Q_x(s),s)] ds - \frac{\pi}{2} \mu(x,t)$  where $U$ is the
  potential which generated the  trajectories. 
\item Define $\tilde{R} (x,t) \ensuremath{:=} P(x) \ensuremath{\lvert {DQ} \rvert}^{-\frac{1}{2}} (x,t)$.
\item Define 
\[  
R(y,t) \ensuremath{:=} 
\sqrt{\sum_{x \in H_{y,t}} \tilde{R}^2 (x,t) + 2 \sum_{\{x_1,
  x_2\}\subset H_{y,t}} \!\tilde{R}(x_1, t)\tilde{R}(x_2,t) 
\cos(S(x_1,t)
  -S(x_2,t))}.
\]


\item Define $u \ensuremath{:=} V  -  \frac{\ensuremath{\hbar}^2}{2}\frac{\ensuremath{\Delta} R}{R}$ 
\item Define $q(x,t) \ensuremath{:=} \gamma_x(t)$ where $\gamma_x$ the solution of
$\ddot{\gamma}_x (s) = -(\nabla u )(\gamma_x(s), s)$ with initial condition 
$\gamma_x(0) = x$, $\dot{\gamma}_x(0) = \nabla T (x)$.  
\end{enumerate}
For an algorithm, fix $T$, $P$, $V$ and define $C(Q,U) \ensuremath{:=}
A2(T,P,V,Q,U)$.
The \ensuremath{{n}^{th}}\ step of A1 is then the same as $C^n(Q_0,U_0)$. Perhaps the
most natural possibility for $Q_0$ in this setting is 
the set of classical trajectories although one could also use a free
motion or a best guess.  


\section{Simple examples}

We shall now work out an  example. 
The potential is $V(x) = \frac{1}{2}x^2$ with initial wave function\
$\ensuremath{\psi} (x,0) \ensuremath{:=} e^{-\frac{1}{2}(x-c)^2+iax}$. Thus, $S(x,0) = ax$ and $R(x,0)
= e^{-\frac{1}{2}(x-c)^2}$. We take $\hbar =1$ throughout this section
and it is more convenient to reindex the iteration  in such a way  that $Q_N$
generates $R_N$ and $S_N$. The actual family of Bohmian trajectories for this
system is  $Q(x,t) = x_0 + c(\cos(t)-1) + a \sin(t) $


Assume that $Q_N(x,t) = x + f_N(t)$ for $f_N \in \ensuremath{C^{\infty}}(\ensuremath{\mathbb{R}^{}})$ such that
$f_N(0) =0$ and $\dot{f_N}(0) = a$. Immediately, we have $Q^{-1}_N (x,t) =
x -f_N(t)$.    
Whether we use A1 or A2, we compute $R_N$ in the same way. We find
that, since $DQ_N^{-1} = 1$,
\[
R_N(x,t) \ensuremath{:=} R(Q^{-1}(x,t), 0) \ensuremath{\lvert {DQ^{-1} (x,t)} \rvert}^{\frac{1}{2}}
= e^{-\frac{1}{2}(x -f_N(t) -c)^2}.
\] The corresponding
quantum potential is 
\[ - \frac{\ensuremath{\Delta} R_N}{2R_N} =
-\frac{1}{2}((x-f_N(t)-c)^2 + 1) =
-\frac{1}{2} (x^2 -2x(f_N(t)+c)
+(f_N(t)+c)^2 +1) .\] 
Thus, we have that the new potential is \[
U_N 
= \frac{1}{2}(2x(f_N(t)+c) - (f_N(t)+c)^2 -1).
\]

For A1, we compute the action integral along $Q_N$ using the potential
$U_N$. 
 We have 
\[
\begin{split} \ensuremath{\mathcal{S}}[Q_N, U_N, S(x,0)](x_0,t)
& = ax_0 + \frac{1}{2}\int_0^t [\dot{f_N}^2(s) - 
  2(x_0+f_N(s))(f_N(s)+c) \\&\quad+ (f_N(s)+c)^2 -1 ]ds \\
& = ax_0 + \frac{1}{2}\int_0^t [-2x_0(f_N(s)+c) + 
 \dot{f_N}^2(s) \\&\quad- f_N(s)^2 +c^2 -1]ds .
\end{split}
\]
To obtain $S_N(x, t)$, we insert $Q_N^{-1}(x,t)$ for $x_0$
in the above formula. To calculate $Q_{N+1}$, we take the gradient of
$S_N$. We compute that  $\nabla S_N(x,t) = a -\int_0^t (f_N(s)+c) ds $


Let $f_{(i), N}$ be the \ensuremath{{i}^{th}} 
antiderivative of $f_N$ such that $f_{(i),N}^{(j)}(0)= 0$ for all $0 \leq j
< i$. 
 We find that
\[
\dot{Q}_{N+1} (x,t) = \nabla S_N (Q_{N+1}(x,t),t) = a - f_{(1),N}(t) - ct .
\] 
The solution to the above equation, with initial condition $Q_{N+1}(x,0) =
x$, is 
\[
Q_{N+1} (x, t) = x+ a t - f_{(2),N}(t) - \frac{1}{2} ct^2  
\] 


As for A2, we do not compute $S_N$ at all. Instead, we compute the
trajectories directly from having the quantum potential.
Indeed, we need to solve
\[
\ddot{Q}_{N+1} (x,t) = -\nabla U_N (Q_{N+1}(x,t), t)
= -(f(t) + c)
\]
with the initial conditions $Q_{N+1}(x,0)= x$ and $\dot{Q}_{N+1}(x,0)
= a$. The solution is 
\[
Q_{N+1} (x,t) = x+ a t - f_{(2),N}(t) - \frac{1}{2} ct^2  
\]

In this example, we have that the two algorithms produce the same
trajectories. Furthermore, the produced solutions are again of the
same form. Starting with $Q_0(x,t) =x_0 +f(t)$ and iterating this
scheme, we find


\[
Q_N (x,t) = x + 
c \sum_{m =1}^{N} \frac{(-1)^m t^{2m}}{(2m)!} +
a \sum_{m =0}^{N} \frac{(-1)^m t^{2m+1}}{(2m+1)!}
+ (-1)^Nf_{(2N)} (t)
.\]
The convergence of the algorithm in terms of trajectories is that of
uniform convergence on $\ensuremath{\mathbb{R}^{}}\times [0,T]$ to the actual Bohmian
trajectories. 


\section{Propagators and the classical algorithm}

We now discuss an application of A2. The idea is to compute the
propagator, or Green's function, for Schr\"odinger's equation 
using this algorithm. If we have a
quadratic potential, then A2 will compute the propagator exactly. The
fixed point trajectories are the classical trajectories. We exemplify
this with the one-dimensional system whose potential is 
 $V(x, t) = \frac{1}{2} a(t) x^2 + b(t) x +c(t)$. 


The Green's function is the integral kernel, $G(x, t; y_0,
0)$, for which 
\[
 \psi (x, t) \ensuremath{:=} \int G(x, t; y_0, 0) \psi(y_0, 0) dy_0
\] is a solution to the Schr\"odinger's equation\ with initial condition $\psi(x,0)$. 
Finding $G(x, t; y_0, 0)$ is equivalent to
solving Schr\"odinger's equation\ with the initial wave function\ being  $\delta (x - y_0)$.  Such a
wave function\ may be thought of as 
representing a swarm of particles located at $y_0$  with a uniform spread of
initial velocities, then we may attempt to use A2 to solve the
problem. Let the uniform distribution be $R(v,0) = \sqrt{\gamma}$
where $\gamma$ is a constant to be discussed later.

Let the initial family of trajectories, which are
now functions of the initial velocity, be of the form
\begin{subequations}
\begin{align} \label{ae:clas}
Q_0(v, t) &= f(t) v + g(t)y_0 +h(t)\\
f(0) = \dot{g}(0) &= h(0) = \dot{h}(0) =  0 \label{ae:clai1}\\
 \dot{f}(0) = & g(0) = 1 \label{ae:clai2}
\end{align}
\end{subequations}
wher $f$, $g$, and $h$ are independent of $y_0$. Then 
\[ Q_0^{-1}(x, t) = \frac{1}{f(t)} (x - y_0 g(t) -h(t) ) .\]
Ignoring the issue for now of when $f(t) = 0$, we find that 
\[
R_0(x,t) \ensuremath{:=} R(Q_0^{-1}(x,t), 0)
\ensuremath{\lvert {\ensuremath{\frac{\partial {Q_0}}{\partial {v}}}} \rvert}^{-\frac{1}{2}}(Q_0^{-1}(x,t)) 
= \sqrt{\frac{\gamma}{f(t)}}.
\] 
We therefore have that the quantum potential vanishes as $R_0$ does
not depend on $x$.  
Thus, $U_0$ is the classical potential and $Q_1$
will be the classical trajectories. 
For a quadratic potential, the classical trajectories are
of the same form as $Q_0$. Thus, the
classical trajectories are the fixed point of the algorithm. 
Therefore, let $f$, $g$, and $h$
be the appropriate functions for the solutions to the classical
system according to \eqref{ae:clas}.
The equations they satisfy, in addition to the initial conditions
\eqref{ae:clai1}, \eqref{ae:clai2} are
\begin{subequations} \label{ae:cla}
\begin{align}
m\ddot{g}(t) &= - a(t)g(t) \\
m\ddot{f}(t) &= - a(t)f(t) \\
m\ddot{h}(t) &= - h(t)a(t) - b(t). 
\end{align} 
\end{subequations}
Note that $h$ vanishes if $b$ vanishes. 


Given $f$, we can immediately write down  $R$  for the Green's
function. 
It is possible for $f$ to have zeros. A zero of $f$, 
which is a very singular caustic, 
  can be dealt with as we suggested before. 
We are essentially  using the square root to see the correction in the phase
appearing as the Maslov index tells us. Since the Maslov index will
increase by 1 through a generic caustic, if we start in the middle of a
caustic, then we take the index to be $\frac{1}{2}$. 
What we need to do is to compute the action $S$; once we have that,
then the solution will be 
\[G(x,t, y_0, 0) = R(t) e^{-i{\pi}{4} - \mu(t)\frac{\pi}{2}}
e^{i\frac{S_{y_{t_{0}}}}{\ensuremath{\hbar}}}. 
\] We
can do this using either the action integral or the fact that $\nabla
S$ is the velocity field, i.e.\ $\nabla S (x,t) = m \dot{Q}(Q^{-1}(x,t),
t)$.  
Define $k_t \ensuremath{:=} y_0 g(t) + h(t)$ and let $f_t$ denote $f(t)$.
 The second method yields 
\[S(x,t) = m(\tilde{S}_{y_0}(t) +
\frac{\dot{f}_t}{f_t}\frac{x^2}{2} 
- \frac{\dot{f}_t}{f_t}(k_t  x
+ (\dot{k}_t x)) \] 
where $\tilde{S}_{y_0}$ is an unknown function for which one needs the first
method to compute. To do so, we write down the action integral, which
is 
\begin{equation*} 
\int_0^t [\frac{m}{2}(\dot{f}_s v + \dot k_s )^2  -(
\frac{1}{2} a_s (f_s v + k_s)^2 +  b_s (f_s v + k_s)
+c_s )] ds.
\end{equation*} 
Replacing $v$ with $\frac{x}{f_t} - \frac{k_t}{f_t}$, 
simplifying, and only retaining the terms not involving $x$, we find 
\begin{equation} \label{ae:int}
\int_0^t [\frac{m}{2}(-\dot{f}_s\frac{k_t}{f_t} +\dot{k}_s)^2
-\frac{1}{2} a(s) (-{f}_s\frac{k_t}{f_t} +{k}_s)^2 +b(s)
 (-{f}_s\frac{k_t}{f_t} +{k}_s) +c(s)] ds
\end{equation}
In summary, we have reduced the computation of the propagator to
computing the above integral and the set of ODEs \eqref{ae:cla}. And,
to simplify even more, the terms not involving $y_0$ are common to all
of the Green's functions and do not have to be computed. Indeed, $c(s)$ is
largely irrelevant and could even be thought of as 
 a free parameter to be chosen to cancel such
terms. 
  
There is also an argument, based on quantum mechanical heuristics, 
which establishes what $\gamma$ should be.
 We want a measure on the velocity space which corresponds
to a delta function in the position space. We therefore take the
Fourier transform of the delta function, finding that
$\frac{1}{\sqrt{2\pi}}e^{iky} \ensuremath{=:} \hat{\ensuremath{\psi}}$ 
with density $\ensuremath{\lvert {\hat{\ensuremath{\psi}}} \rvert}^2 =
\frac{1}{2\pi}$. But this is momentum space and the relation to the
velocity space is $k = \frac{mv}{\ensuremath{\hbar}}$. Taking the differential of that
relation, we find that $\frac{dk}{2\pi} = \frac{m dv}{2\pi\ensuremath{\hbar}}$
implying that $\gamma = \frac{m}{2\pi \ensuremath{\hbar}}$. 

We finish by explicitly computing the propagators for the
free potential and the harmonic oscillator. The free potential, $V(x)
=0$, has $f(t) = t$, $g(t) = 1$, and $h(t) = 0$.  We
find that the propagator is 
\[
 G(x,t, y_0, 0) = 
\sqrt{\frac{\gamma}{t}} e^{-i\frac{\pi}{4}} e^{i \frac{m}{\ensuremath{\hbar}} 
\frac{(x-y_0)^2}{2t}}
\]
for $t > 0$. For $t<0$ multiply the above by $e^{i\frac{\pi}{2}}$ as
the Maslov index indicates. It is a fact that 
\[
\delta (x-y_0) = \lim_{a \to +\infty} \frac{a}{\sqrt{\pi}} e^{i(a^2
(x-y_0)^2 - \frac{\pi}{4})} . 
\]
Thus, this confirms that the initial condition is satisfied, that the
Maslov index is contributing the correct factor, and that
 $\gamma$ is $ \frac{m}{2\ensuremath{\hbar}
\pi}$.
   
 
As for the harmonic oscillator, the following comes out. Let the
potential be $V(x) = \frac{m\omega^2}{2} x^2$. Then $f(t) =
\frac{1}{\omega} \sin(\omega t)$, $g(t) = \cos(\omega t)$, and $h(t) = 0$.  
The integral \eqref{ae:int} becomes 
\[
\begin{split}
\tilde {S}_{y_0}(t) =& 
\int_0^t [\frac{m}{2}(\cos(\omega s) \frac{\omega y_0 \cos(\omega
t)}{\sin(\omega t)} + y_0 \omega \sin(\omega s) )^2
\\& -\frac{m \omega^2}{2}  (-\frac{\sin(\omega s)}{\omega} 
 \frac{\omega y_0 \cos(\omega
t)}{\sin(\omega t)}+y_0 \cos(\omega
s))^2 ] ds \\
=&\frac{m \omega ^2 y_0}{2} \int_0^t[(\cos^2 (\omega s) - \sin^2(\omega
s))  \cot(\omega t)^2 
\\& \hspace{.5in} +4 (\cos(\omega s) \sin(\omega s)) \cot(\omega t)   
\\& \hspace{.5in} -( \cos^2 (\omega s) - \sin^2(\omega s) )]ds \\ 
\\=&  \frac{m \omega ^2 y_0}{2} (\frac{1}{\omega}(\sin (\omega
t)\cos(\omega t)  (\cot(\omega t)^2 -1) 
+2 (\frac{1}{\omega} \sin^2(\omega t)) \cot(\omega t)   
\\=& 
\frac{m \omega y_0}{2} \cot(\omega t) 
\end{split}
\]
Thus, when we put it all together, we find 
\[
G(x,t, y, 0) = 
\sqrt{\ensuremath{\lvert {\frac{\gamma \omega}{\sin(\omega t)}} \rvert}} e^{-i(\frac{\pi}{4}
+ \mu(t) \frac{\pi}{2})}  
e^{i \frac{m \omega}{2\ensuremath{\hbar}} 
(x^2+y^2)\cot(\omega t) - 2 xy \csc(\omega t)}.
\]
The Maslov index, $\mu(t)$, starts at $0$ and is incremented by one at
the end of every half period; this incrementation corresponds to when
all of the trajectories come together and 
form a very singular caustic. 
The Maslov index produces the correct
phase as the trajectories move through the singularity. 
Explicitly, we have $\mu(t) = k $ for $t \in [k\frac{\pi}{\omega},
(k+1) \frac{\pi}{\omega})$. Computing the limits as $t$
approaches a singular point, we find that 
\[
G(x, k\frac{\pi}{\omega}, y, 0) = e^{i \frac{\pi}{2} k} 
\begin{cases}
\delta (x-y) &\text{if $k$ is even} \\
\delta (x+y) &\text{if $k$ is odd}  
\end{cases}.
\]
This is a complete description of the propagator for the harmonic
oscillator.


\section{Some calculations involving the chain rule}

We wish to provide a succinct formula for computing the quantities in
these algorithms. The basic quantity to compute is the quantum force
and this involves three derivatives of an object created by
composition.  These are all simple calculations using the chain rule. We shall
use the notations of Section \ref{as:short}


We start with $R$. This is the object for which much computation must
be done. First, we define $\ensuremath{\tilde R} (x,t) \ensuremath{:=} P(x)
\ensuremath{\lvert {DQ} \rvert}^{-\frac{1}{2}}$ and then $R \ensuremath{:=} \ensuremath{\tilde R} (H)$. We will have to compute
up to three derivatives of $R$. Using the chain rule,
we find that
\begin{align*}
DR (v_1) =& D\ensuremath{\tilde R}(DH(v_1)) \\
D^2 R (v_1 \otimes v_2) =& 
D^2\ensuremath{\tilde R} (DH(v_1)\otimes DH(v_2)) +D\ensuremath{\tilde R}(D^2H(v_1 \otimes v_2) )  
\\
\begin{split}
D^3 R (v_1 \otimes v_2 \otimes v_3)  =& 
D^3\ensuremath{\tilde R} (DH (v_1) \otimes DH (v_2 ) \otimes DH (v_3) )  \\
& +D^2\ensuremath{\tilde R} \{
(D^2H(v_1\otimes v_2) \otimes DH (v_3)) \\& 
\hspace{.43in}+(D^2H(v_1\otimes v_3) \otimes DH (v_2)) 
\\&\hspace{.43in}+
D_3H (v_1) \otimes (D^2H(v_2\otimes v_3)) \} \\
& + D\ensuremath{\tilde R}(D^3H (v_1\otimes v_2 \otimes v_3)) .
\end{split}
\end{align*}

In our equations, we need to compute the quantum force which is the gradient of
the Laplacian of $R$ divided by $R$. 


We have that 
\[
D( \frac{\ensuremath{\Delta} R}{R})(v) = \frac{D^3 R (G\otimes v)}{R} - \frac{D^2 R (G)
DR(v)}{R^2} 
\]
where $G \in TM^{\otimes 2}$ is the dual of $g$; in coordinates, $G =
\sum_{(i,j) \in \ensuremath{\mathfrak{n}}^2} G^{ij} v_i \otimes v_j $ and $G^{ij}$ are the
components of the inverse matrix of $g_{ij}$.

The derivatives of \ensuremath{\tilde R}\ can in turn be computed using the product rule,
the chain rule and the differentiation formulas for the determinant
developed in Appendix \ref{a:app7}. 


Both for implementing
 the algorithms and for an application of A2 described in the
next section, we need to compute $D^i_{Q(x,t)} [Q^{-1}]$ in terms  of the
derivatives of $DQ$. Yet again we use the chain rule.  We need it up
to the third order. Let $y \ensuremath{:=} Q(x)$; alternatively, define $x \ensuremath{:=}
Q^{-1} (y)$. We also define $\ensuremath{B} \ensuremath{:=} (D_x [Q])^{-1}$. For
$i=1,2,3$, let  $v_i
\in T_yM$ be arbitrary. 
 We then find that
\begin{align}
Q^{-1} (Q(x)) =& \ensuremath{\text{Id}} \\
D_{y} [Q^{-1}] (v_1) =& B (v_1) \\
D_{y}^2 [Q^{-1}](v_1 \otimes v_2)
 =& -B (D^2_x [Q] (B (v_1) \otimes B (v_2))) \\
\begin{split}
D_{y}^3 [Q^{-1}] (v_1\otimes v_2\otimes v_3)  =& 
-B (D^3_x [Q] \{B(v_1)\otimes B(v_2) \otimes
B(v_3) \} ) \\ 
& + B (D^2_x [Q] (B(  
D_x^2 [Q] ( B (v_1) \otimes B (v_2)))  \otimes B(v_3)   \\
& \hspace{.68in} + B(D_x^2 [Q] ( B (v_1) \otimes B (v_3)))  \otimes
B(v_2)  \\ 
& \hspace{.68in} + B( v_1) \otimes B(D_x^2 [Q] ( B (v_2) \otimes B (v_3) ))) 
\end{split}
\end{align}
These are routine calculations. For example, to
compute the second derivative, the chain rule leads to $D^2_Q
[Q^{-1}](D_x [Q] \otimes D_x[Q]) + D_Q [Q^{-1}] (D^2_x Q) = 0$. One then
solves for the second derivative of the inverse and then composes
with the inverse of the first derivative using 
the equation
for the first derivative to substitute in $B$.

In total, we can see what we need. Before starting the algorithms, we
compute the first derivative of $T$, i.e.\ 
the 1-form dual to the initial Bohmian velocity field. We also compute three
derivatives of $P$. The rest of the information is mainly computed
from $Q$. 
We need to compute four derivatives of $Q$ although 
the fourth derivative appears only in the determinant derivative. 
To finish the computations, $Q^{-1}$ and $DQ^{-1}$ need to be
computed. We then simply plug the objects into the correct places.  


\section{The Bohmian velocity partial differential equation}

From the previous section, we discussed some ways in which the
essential computations can be organized. We apply these ideas here, at
the fixed point, to derive a PDE for the family of Bohmian trajectories. 
One advantage of this is that it is not necessary to compute the
inverse of a given family of trajectories; the disadvantage is that
we are again facing a PDE. 

The idea is that one writes down the family of ODE's generated by the
algorithm  A2. One then assumes that the fixed point has been
achieved. The ODE for A2 is 
\[
g(\ddot{q} (x,t), ) = -dV(q(x,t)) + QF(Q^{-1}(q(x,t),t),t)
\]
where the quantum force, $QF$, is thought of as being a function of
the initial $x_0$. When we evaluate along the fixed point, we
have that $Q^{-1}$ and $Q$ cancel, just leaving the quantum force evaluated
at $x$. 

We have the following PDE for $Q$, the family of Bohmian trajectories
for a given $\ensuremath{\psi}(x,0)$ and potential $V$. We write it in terms of
1-forms acting on a vector $v$. 
$C$ is the determinant of $DQ$, $B$ is $(D[Q])^{-1}$. We
sometimes employ subscripts to denote which objects are being contracted.  
A rough schematic of the quantum force expression is quite helpful in
understanding the PDE; the schematic is
\[
\frac{1}{\ensuremath{\tilde R}}\{D^3 [Rq] (\cdot) +D^2[\ensuremath{\tilde R}] (\cdot) +D[\ensuremath{\tilde R}] (\cdot) \}
- \frac{1}{\ensuremath{\tilde R}^2} \{D^2[\ensuremath{\tilde R}](\cdot) +D[\ensuremath{\tilde R}](\cdot)  \}\{D[\ensuremath{\tilde R}](\cdot) \}.
\]
One task is to compute, somewhat explicitly, the derivatives of
\ensuremath{\tilde R}. The other task is to compute the various different tensors that
are composed into the $\cdot$ slots. It can get quite messy.

\[
\begin{split}
g(\ddot{Q}(x,t),v ) = & -dV_{Q} (v) + \frac{\ensuremath{\hbar}^2}{2}(
\frac{1}{P} \{
\\&
 (D^3 P - \frac{15}{8} P C^{-3} D^3C 
\\ &
- \frac{1}{2} C^{-1} D^2 P_{12} \otimes DC_3 
- \frac{1}{2} C^{-1} D^2 P_{13} \otimes DC_2 
- \frac{1}{2} C^{-1} DC_1 \otimes D^2P_{23} 
\\ & 
+ \frac{3}{4}  C^{-2} D^2C_{12} \otimes DP_3 
+ \frac{3}{4}  C^{-2} D^2C_{13} \otimes DP_2 
+ \frac{3}{4}  C^{-2} DP_1 \otimes D^2C_{23} )
\\&\qquad \circ ( \{(B\otimes B)\circ G\} \otimes B(v) )
\\ & 
+( D^2 P 
+ \frac{3}{4} P C^{-2} D^2C
\\ &
- \frac{1}{2} C^{-1} DP \otimes DC 
- \frac{1}{2} C^{-1} DC \otimes DP )
\\&\qquad \circ 
( -\{ B \circ (D^2 [Q])\circ (B
\otimes B)   \circ G \} \otimes B(v)  
\\ &\qquad \quad
- \{B \circ (D^2 [Q])\circ (B (G_1)
\otimes B( v))\} \otimes B(G_2)  
\\ & \qquad \quad
- B(G_1) \otimes \{B \circ (D^2 [Q])\circ (B (G_2)
\otimes B(v)) \}  )
\\ &
+(C^{-\frac{1}{2}} DP -
\frac{1}{2} P C^{-1} DC)
\\&
\qquad \circ (
 -B \circ (D^3 [Q]) \circ \{(B \otimes B) \circ G\} \otimes
B(v)
\\ & 
\qquad \quad + B \circ D^2 [Q] \circ B\otimes B \circ (
\{ D^2 [Q] \circ ( B \otimes B) \circ G\} \otimes v   \\
& \hspace{2.05in} + \{D^2 [Q] \circ ( B(G_1) \otimes B (v) )\}  \otimes
G_2  
\\& \hspace{2.05in} +  G_1 \otimes \{ D^2 [Q] ( B (G_2) \otimes B
(v)\}  ) ) )\}
\\ &
- \frac{1}{P^2}   
\{
( D^2 P 
- \frac{1}{2} C^{-1} DP \otimes DC 
- \frac{1}{2} C^{-1} DC \otimes DP 
+ \frac{3}{4} P C^{-2} D^2C )
\\&\qquad\circ(B\otimes B)\circ G
\\ &
+ ( DP -\frac{1}{2} P C^{-1} DC) \circ (-B) \circ (D^2 [Q])\circ (B
\otimes B) \circ G  \} 
\\ &
\{  DP -
\frac{1}{2} P C^{-1} DC) (B (v)) \}   
\end{split}   
\]

with the initial conditions
\begin{align*}
Q(x,0) =& x \\
\dot{Q}(x,0) =& \nabla T (x).  
\end{align*}
For the algorithms, the only change in the quantum force 
is the evaluation which takes place at 
$Q^{-1} (q(x,t))$ instead of at $x$;
 appearances of $Q$ outside the quantum force term
are replaced with $q$.  

There is not much to say about this PDE. Notice that if $C$ is
constant, then many of the terms vanish. The equation, however is
still rather unapproachable.


A key simplification occurs when  $D^2Q = 0$ everywhere which
  geometrically is when the mapping $Q$ preserves geodesics as
  proven in Chapter \ref{c:taylor}. 
The PDE becomes the relatively simple
\[
\begin{split}
g(\ddot{Q}(x,t),v ) = & -dV_{Q} (v) + \frac{\ensuremath{\hbar}^2}{2}(
\frac{1}{P} 
 D^3 P \circ ( \{(B\otimes B)\circ G\} \otimes B(v) )
\\ &
- \frac{1}{P^2}   
\{ D^2 P 
\circ(B\otimes B)\circ G
 \} 
\{ DP  (B (v)) \}.   
\end{split}   
\]


Notice that $C$ has disappeared entirely from the equation. 


\subsection{The one-dimensional case}

If $M$ is just \ensuremath{\mathbb{R}^{}}, then $G$ becomes $\frac{1}{m}$, the
derivatives become numbers, and we use primes to denote
spatial differentiation. The full PDE 
simplifies slightly,  but it is still brutal. 

If $Q'' = 0$ identically, then we can solve the equation. We have
\[
m \ddot{Q}(x,t)  =  - V_{Q}'  + \frac{\ensuremath{\hbar}^2}{2}(
\frac{P'''}{m P (Q')^{3} } 
- \frac{P''P'}{m P^2 (Q')^3}.)   
\]


Furthermore, $Q'' =0$ implies that $Q(x,t) = f(t) x + k(t)$. 

Let $H_n (x)$ denote the \ensuremath{{n}^{th}}\ Hermite polynomial. 
We claim that if $V(x) = \frac{m}{2} \omega^2 x^2 + m \gamma x
$ and 
\[\ensuremath{\psi}(x,0) =
H_n(\sqrt{md/\ensuremath{\hbar}}(x-c))  e^{-\frac{md}{2\ensuremath{\hbar}} (x-c)^2} e^{i\frac{m}{\ensuremath{\hbar}}
(\frac{a}{2} (x-c)^2 +b(x-c))},\] ($d>0$), 
then the trajectory solutions are of the above affine form and this PDE can be
solved quite explicitly; it actually becomes two decoupled ODEs.  We shall
use  that $y(x) \ensuremath{:=} H_n(x)e^{-\frac{x^2}{2}}$ satisfies $y'' = (x^2- 2n
-1)y$; this comes from the Schr\"odinger\ equation for the harmonic oscillator. 
For us, the useful imphlication of the Hermite differential equation is 
\[
y y''' -y'' y' = 2x y^2 +(x^2 -2n-1)y'y -(x^2 -2n-1)yy' = 2xy^2
\]
When we use it, $x$ will be $\sqrt{md/\ensuremath{\hbar}}(x-c)$ and we shall
also be using the chain rule which produces a multiplicative factor of 
$\sqrt{md/\ensuremath{\hbar}}^{3}$.

It is convenient to take $Q(x,t) = f(t) (x-c) + k(t)$. Then we have
that 
\[ 
\frac{PP''' -P''P'}{P^2} =2(\sqrt{md/\hbar}
(x-c)\sqrt{md/\hbar}^3 
\] 
The PDE becomes
\[
\begin{split}
m(\ddot{f}(t) (x-c)  + \ddot{k}(t)) =& \frac{\ensuremath{\hbar}^2}{2m f(t)^3}
(2(md/\ensuremath{\hbar})^2(x-c)) 
\\& - m\omega^2 (f(t) (x-c)  + k(t)) - m\gamma 
\\
= & \frac{m d^2}{f(t)^3} (x-c)-  
m\omega^2 f(t) (x-c) - m \omega^2 k(t) -m\gamma.  
\end{split}
\]
This decouples into the two independent ODEs
\begin{align}
\ddot{f} =& \frac{d^2}{f^3} - \omega^2 f \label{ae:fode} \\
\ddot{k} =& -\omega^2 k - \gamma \label{ae:kode}
\end{align}
with the initial conditions
\[
f_0 = 1, \dot{f}_0  = a , k_0 = c, \dot{k}_0 = b . 
\]

The ODE \eqref{ae:kode} is an easy equation to solve; its solution is 
\[
k(t) = (c + \frac{\gamma}{\omega^2})\cos(\omega t) +
\frac{b}{\omega} \sin(\omega t) - 
\frac{\gamma}{\omega^2} .
\]
The other ODE is not as easy to solve. The inspiration is
that \eqref{ae:fode} is actually the ODE that the radius function
satisfies in a 2-dimensional, classical, harmonic oscillator
system. More precisely, let $u$ and $v$ satisfy 
\begin{align*} 
\ddot{u} = - \omega^2 u,\quad & \ddot{v} = -\omega^2v \\
u(0) = 1, \dot{u}(0) = a ,\quad& v(0) = 0, \dot{v}(0) = d.
\end{align*}
Define $r \ensuremath{:=} \sqrt{u^2 +v^2}$. Then, by differentiating and using
the conservation of $J^2 = (\dot{u}v -\dot{v}u)^2 = d^2$, 
we find that $r$
satisfies 
\[
\ddot{r} = \frac{d}{r^3} - \omega^2 r 
\]
with the initial conditions
\[
r(0) = 1, \dot{r}(0) =  a.
\]
Thus, $r$ and $f$ satisfy the same equations implying that they are
the same function. But we know how to
solve the harmonic oscillator. Doing so and finding the radial
function, we have
\[
f(t) = r(t) = \sqrt{\cos^2 (\omega^2 t) + 2 \frac{a}{\omega} \cos(\omega
t) \sin (\omega t) + \frac{a^2 + d^2}{\omega^2} \sin^2 (\omega t)}.
\]
For both $f$ and $k$, if we take the limit as $\omega$ goes to zero,
then the correct solution for $\omega =0$ is the limiting solution. 
 

We have just computed the Bohmian trajectories. An immediate
consequence is that we have the
amplitude. Namely, 
\[
\begin{split}
R(x,t) =& P( \frac{x -k(t)}{f(t)} + c)\ensuremath{\lvert {f(t)} \rvert}^{-\frac{1}{2}}
\\ &
= \frac{1}{\sqrt{f(t)}}
H_n(\sqrt{md/\ensuremath{\hbar}}(\frac{x -k(t)}{f(t)})) e^{-\frac{md}{2\ensuremath{\hbar}}
(\frac{x -k(t)}{f(t)})^2}.
\end{split}
\] 
The phase is 
\[
S(x,t) = \frac{m}{\ensuremath{\hbar}} (\frac{1}{2} \dot{f}(t) (\frac{x -k(t)}{f(t)})^2 +
\dot{k}(t) (\frac{x-k(t)}{f(t)}) +h(t)) 
\]
where $h(t)$ is a function for which one needs to do the action
integral in order to determine it. We shall not compute $h(t)$.

\clearpage
\bibliographystyle{abbrv}
\bibliography{references}
\end{document}
